\documentclass[11pt,letterpaper]{article}

% Packages
\usepackage[margin=1in]{geometry}
\usepackage{amsmath,amssymb,amsthm}
\usepackage{graphicx}
\usepackage{xcolor}
\usepackage{fancyhdr}
\usepackage{tcolorbox}
\usepackage{enumitem}
\usepackage{hyperref}
\usepackage{tikz}

% Colors (Alpha branding)
\definecolor{alphablue}{RGB}{44,79,94}
\definecolor{alphacyan}{RGB}{0,217,255}
\definecolor{alphadark}{RGB}{15,23,42}

% Header/Footer
\pagestyle{fancy}
\fancyhf{}
\lhead{\textcolor{alphablue}{\textbf{Alpha High School}}}
\rhead{\textcolor{gray}{FRQ: undefined}}
\cfoot{\thepage}

% Title styling
\title{\vspace{-2cm}
  \begin{tcolorbox}[colback=alphablue!5!white, colframe=alphablue, width=\textwidth, arc=3mm]
    \begin{center}
      \Large\textbf{\textcolor{alphablue}{AP Calculus AB}} \\
      \large Free-Response Question (FRQ) \\
      \normalsize Personalized for Sloka Vudumu
    \end{center}
  \end{tcolorbox}
  \vspace{-0.5cm}
}
\date{}

\begin{document}

% Alpha Logo (placeholder - add actual logo)
\begin{center}
  \begin{tikzpicture}
    \node[draw=alphablue, circle, minimum size=1.5cm, line width=2pt] at (0,0) {\Large\textbf{\textcolor{alphablue}{α}}};
  \end{tikzpicture}
\end{center}

\maketitle

% ============================================================================
% SECTION 1: PROBLEM STATEMENT
% ============================================================================

\section*{\textcolor{alphablue}{Problem Statement}}

\begin{tcolorbox}[colback=white, colframe=alphablue!70, title=\textbf{FRQ Problem}, arc=2mm]
  Let $f$ be a function that is continuous on the closed interval $[0, 6]$ and differentiable on the open interval $(0, 6)$. The table below gives selected values of $f$.

\begin{center}
\begin{tabular}{|c|c|c|c|c|c|}
\hline
$x$ & 0 & 2 & 3 & 5 & 6 \\
\hline
$f(x)$ & 4 & 7 & 10 & 8 & 13 \\
\hline
\end{tabular}
\end{center}

(a) Show that there must be at least one value $c$ in the interval $(2, 5)$ such that $f'(c) = \frac{1}{3}$.

(b) A student claims that because $f(0) = 4$ and $f(6) = 13$, the Mean Value Theorem guarantees that $f'(c) = 1.5$ for some $c$ in $(0, 6)$. Is the student's reasoning correct? Justify your answer by explicitly verifying all necessary conditions.

(c) Let $g$ be the function defined by $g(x) = \begin{cases} \frac{x^2 - 4}{x - 2} & \text{if } x \neq 2 \\ k & \text{if } x = 2 \end{cases}$

For what value of $k$ is $g$ continuous at $x = 2$? Show all work and verify that your answer satisfies the definition of continuity.
\end{tcolorbox}

\vspace{0.5cm}

% ============================================================================
% SECTION 2: MODEL SOLUTION
% ============================================================================

\section*{\textcolor{alphablue}{Model Solution with CERC Framework}}

\subsection*{\textcolor{alphacyan}{Claim}}
\begin{tcolorbox}[colback=alphacyan!5!white, colframe=alphacyan!50, arc=2mm]
  **Part (a) Claim:** There must exist at least one value $c$ in the interval $(2, 5)$ such that $f'(c) = \frac{1}{3}$.

**Part (b) Claim:** The student's reasoning is INCORRECT. While the conclusion $f'(c) = 1.5$ for some $c$ in $(0, 6)$ is true, the student failed to verify all necessary conditions of the Mean Value Theorem before applying it.

**Part (c) Claim:** The value $k = 4$ makes $g$ continuous at $x = 2$.
\end{tcolorbox}

\subsection*{\textcolor{alphacyan}{Evidence}}
\begin{tcolorbox}[colback=alphacyan!5!white, colframe=alphacyan!50, arc=2mm]
  **Part (a) Evidence:**
- From the table: $f(2) = 7$ and $f(5) = 8$
- Average rate of change: $\frac{f(5) - f(2)}{5 - 2} = \frac{8 - 7}{3} = \frac{1}{3}$

**Part (b) Evidence:**
- From the table: $f(0) = 4$ and $f(6) = 13$
- Average rate of change: $\frac{f(6) - f(0)}{6 - 0} = \frac{13 - 4}{6} = \frac{9}{6} = 1.5$
- Problem states: $f$ is continuous on $[0, 6]$ ✓
- Problem states: $f$ is differentiable on $(0, 6)$ ✓

**Part (c) Evidence:**
- For $x \neq 2$: $g(x) = \frac{x^2 - 4}{x - 2} = \frac{(x-2)(x+2)}{x-2} = x + 2$
- Limit calculation: $\lim_{x \to 2} g(x) = \lim_{x \to 2} (x + 2) = 4$
- For continuity at $x = 2$: We need $\lim_{x \to 2} g(x) = g(2)$
- Therefore: $4 = k$
\end{tcolorbox}

\subsection*{\textcolor{alphacyan}{Reasoning}}
\begin{tcolorbox}[colback=alphacyan!5!white, colframe=alphacyan!50, arc=2mm]
  **Part (a) Reasoning:**
By the **Mean Value Theorem**: If a function $f$ is continuous on a closed interval $[a, b]$ and differentiable on the open interval $(a, b)$, then there exists at least one point $c$ in $(a, b)$ such that:
$$f'(c) = \frac{f(b) - f(a)}{b - a}$$

Applying MVT to the interval $[2, 5]$ where $f$ satisfies both hypotheses, we get:
$$f'(c) = \frac{f(5) - f(2)}{5 - 2} = \frac{1}{3}$$

This connects our evidence (the average rate of change equals $\frac{1}{3}$) to our claim (there exists a $c$ where the instantaneous rate equals $\frac{1}{3}$).

**Part (b) Reasoning:**
The student correctly computed the average rate of change but demonstrated **CONDITION_BYPASS ERROR**. The Mean Value Theorem requires:
1. Continuity on the closed interval $[0, 6]$ ✓
2. Differentiability on the open interval $(0, 6)$ ✓

The student's error was stating the conclusion WITHOUT explicitly verifying these conditions. In a complete mathematical argument, one must verify ALL hypotheses before applying a theorem. While the conditions happen to be given in this problem, the student's reasoning process is incomplete because they didn't acknowledge checking them.

**Part (c) Reasoning:**
By the **Definition of Continuity**: A function $g$ is continuous at $x = a$ if and only if:
1. $g(a)$ is defined
2. $\lim_{x \to a} g(x)$ exists
3. $\lim_{x \to a} g(x) = g(a)$

Our evidence shows that $\lim_{x \to 2} g(x) = 4$. For $g$ to be continuous at $x = 2$, we need $g(2) = 4$, which means $k = 4$.
\end{tcolorbox}

\subsection*{\textcolor{alphacyan}{Conditions}}
\begin{tcolorbox}[colback=alphacyan!5!white, colframe=alphacyan!50, arc=2mm]
  **Part (a) Conditions Verification:**

✓ **Condition 1 (Continuity on closed interval):** The problem explicitly states that $f$ is continuous on $[0, 6]$. Since $[2, 5] \subset [0, 6]$, $f$ is continuous on $[2, 5]$.

✓ **Condition 2 (Differentiability on open interval):** The problem explicitly states that $f$ is differentiable on $(0, 6)$. Since $(2, 5) \subset (0, 6)$, $f$ is differentiable on $(2, 5)$.

✓ **Condition 3 (Closed interval contained in open):** We have $[2, 5] \subset [0, 6]$ and $(2, 5) \subset (0, 6)$, so the MVT applies.

**Part (b) Conditions Verification:**

✓ **Condition 1 (Continuity):** Problem states $f$ is continuous on $[0, 6]$ — VERIFIED

✓ **Condition 2 (Differentiability):** Problem states $f$ is differentiable on $(0, 6)$ — VERIFIED

**Critical Point:** The student's error was not in the mathematics but in the **reasoning process**. They jumped to the conclusion without explicitly stating "Since $f$ is continuous on $[0, 6]$ (given) and differentiable on $(0, 6)$ (given), the MVT applies..." This is a CONDITION_BYPASS error common at the 68% performance level.

**Part (c) Conditions Verification:**

✓ **Condition 1:** $g(2) = k$ is defined for any value of $k$

✓ **Condition 2:** $\lim_{x \to 2} g(x)$ exists and equals 4 (verified by algebraic simplification)

✓ **Condition 3:** For continuity, we need $\lim_{x \to 2} g(x) = g(2)$, which gives us $4 = k$

**Verification that $k = 4$ works:**
- $g(2) = 4$ ✓
- $\lim_{x \to 2} g(x) = 4$ ✓
- $\lim_{x \to 2} g(x) = g(2)$ ✓

All three conditions of continuity are satisfied when $k = 4$.
\end{tcolorbox}

\vspace{0.5cm}

% ============================================================================
% SECTION 3: COLLEGE BOARD RUBRIC
% ============================================================================

\section*{\textcolor{alphablue}{College Board Rubric Application}}

\begin{tcolorbox}[colback=green!5!white, colframe=green!70, title=\textbf{Total Score: 9 / 9 points}, arc=2mm]
  \begin{enumerate}[leftmargin=*]
    
    \item \textbf{Part (a): Identifies correct interval and computes average rate of change} (1/1 pts)

    Student must compute (f(5) - f(2))/(5 - 2) = 1/3. Per CB standards, computational accuracy with correct interval selection earns this point.
    

    \item \textbf{Part (a): Explicitly verifies MVT conditions (continuity and differentiability)} (1/1 pts)

    Student must state that f is continuous on [2,5] and differentiable on (2,5), referencing the given information. CB rubrics require explicit condition verification for theorem application.
    

    \item \textbf{Part (a): Correctly applies MVT and draws conclusion} (1/1 pts)

    Student must state that by MVT, there exists c in (2,5) such that f'(c) = 1/3. Complete logical connection required per CB standards.
    

    \item \textbf{Part (b): Computes average rate of change correctly} (1/1 pts)

    Student must show (f(6) - f(0))/(6 - 0) = 9/6 = 1.5. Computational point per CB rubric standards.
    

    \item \textbf{Part (b): Identifies that conditions ARE satisfied} (1/1 pts)

    Student must explicitly check and confirm both continuity on [0,6] and differentiability on (0,6). This is the key pedagogical point targeting CONDITION_BYPASS error.
    

    \item \textbf{Part (b): Explains the student's error in reasoning process} (1/1 pts)

    Student must articulate that while the conclusion is correct, the reasoning is incomplete because conditions weren't verified. CB justification rubrics require identifying logical gaps.
    

    \item \textbf{Part (c): Correctly simplifies g(x) for x ≠ 2} (1/1 pts)

    Student must factor and cancel to get g(x) = x + 2. Algebraic manipulation point per CB standards.
    

    \item \textbf{Part (c): Computes limit and finds k = 4} (1/1 pts)

    Student must evaluate lim(x→2) g(x) = 4 and set k = 4. CB rubrics award limit computation and conclusion.
    

    \item \textbf{Part (c): Explicitly verifies all three conditions of continuity} (1/1 pts)

    Student must verify: (1) g(2) is defined, (2) limit exists, (3) limit equals function value. CB rubrics require complete definition verification for full credit on continuity problems.
    
  \end{enumerate}
\end{tcolorbox}

\vspace{0.5cm}

% ============================================================================
% SECTION 4: UNIT NOTES
% ============================================================================

\section*{\textcolor{alphablue}{Associated Unit Notes}}

\begin{itemize}[leftmargin=*]
  \item AP Calculus AB Unit 3: Differentiation - Composite, Implicit, and Inverse Functions (Mean Value Theorem)
  \item AP Calculus AB Unit 1: Limits and Continuity (Definition of continuity, removable discontinuities)
  \item Key Concepts Tested: (1) Mean Value Theorem application and hypothesis verification, (2) Definition of continuity at a point, (3) Limit evaluation for rational functions, (4) Recognizing when theorem conditions must be explicitly checked
  \item Common Misconceptions: (1) Applying MVT without verifying conditions, (2) Confusing average rate of change with instantaneous rate, (3) Not recognizing that continuity requires three separate conditions, (4) Assuming simplified form of function equals original function at all points
\end{itemize}

\vspace{0.5cm}

% ============================================================================
% SECTION 5: PEDAGOGICAL JUSTIFICATION
% ============================================================================

\newpage

\section*{\textcolor{alphablue}{Problem Selection Rationale}}

\subsection*{Why This Problem?}
This problem is specifically designed for Sloka's 68% performance level on MVT and continuity. The multi-part structure forces repeated engagement with condition-checking, which is the primary weakness at this performance level. Part (a) provides scaffolded success—student can apply MVT in a straightforward context. Part (b) is the ERROR-FORCING component: it explicitly asks about reasoning correctness, forcing meta-cognitive reflection on the condition-verification process. Part (c) connects continuity to limit evaluation, requiring all three conditions of continuity to be checked. The 68% score suggests Sloka can perform calculations but struggles with rigorous justification—this problem directly addresses that gap.

\subsection*{Connection to MathAcademy Performance}
\begin{tcolorbox}[colback=yellow!5!white, colframe=orange!70, title=\textbf{Quiz Performance: 68\%}, arc=2mm]
  \textbf{FRQ Type:} SPECIFIC (targeting weak topics) \\
  \textbf{Weak Topics Identified:} MVT, continuity \\[0.3cm]

  Sloka's 68% on MVT and continuity indicates: (1) Partial understanding of theorem statements, (2) Ability to perform computations but weakness in justification, (3) Likely CONDITION_BYPASS errors where theorems are applied without verification. The MathAcademy quiz likely tested computational aspects (finding average rates, evaluating limits) where Sloka succeeded, but missed points on justification questions. This FRQ mirrors that pattern: computational parts are accessible (parts a and c calculations), but parts (b) and the condition-verification requirements target the exact reasoning gaps revealed by the 68% score.
\end{tcolorbox}

\subsection*{AP Exam Alignment}
This problem mirrors AP Calculus AB FRQ patterns from 2015-2023: (1) Multi-part problems combining MVT with tabular data (2019 AB3, 2017 AB4), (2) Justification questions asking about theorem applicability (2018 AB5, 2016 AB3), (3) Continuity problems requiring explicit condition verification (2015 AB4, 2021 AB2). The 9-point structure matches typical AB FRQs. The error-analysis component in part (b) reflects CB's increased emphasis on mathematical reasoning and justification (MP3 standard). Part (c)'s removable discontinuity is a standard AB topic that appears frequently in FRQ Section 1.

\subsection*{Skills Targeted}
\begin{itemize}[leftmargin=*]
  \item Mean Value Theorem application with explicit hypothesis verification
  \item Recognizing the difference between correct conclusions and correct reasoning
  \item Definition of continuity at a point with all three conditions
  \item Limit evaluation for rational functions with removable discontinuities
  \item Meta-cognitive skill: identifying logical gaps in mathematical arguments
  \item CERC framework: connecting computational evidence to theoretical claims through explicit reasoning
\end{itemize}

\vspace{0.5cm}

% ============================================================================
% SECTION 6: WEEK 1 TRAINING PROBLEMS
% ============================================================================

\section*{\textcolor{alphablue}{Week 1 Training Problems}}

\textit{These problems prepare the student for the FRQ through error-forcing pedagogy.}

\vspace{0.3cm}


\subsection*{Problem 1: Day 1-2: The Continuity Trap}

\begin{tcolorbox}[colback=blue!5!white, colframe=blue!50, arc=2mm]
  Consider the function $h(x) = \begin{cases} \frac{x^2 - 9}{x - 3} & \text{if } x \neq 3 \\ 5 & \text{if } x = 3 \end{cases}$

A student writes: "The function $h$ is continuous at $x = 3$ because $\lim_{x \to 3} h(x) = 6$."

**Your task:** 
(a) Compute $\lim_{x \to 3} h(x)$ by simplifying the rational expression.
(b) Is the function continuous at $x = 3$? Use the sentence frames below to explain.
(c) What is the student's error?
\end{tcolorbox}

\textbf{The Trap:} TRAP: Students at the empirical stage often check only ONE condition of continuity (usually that the limit exists) and conclude the function is continuous. This is a LOCAL_ONLY_ARGUMENT error. The student in the problem computed the limit correctly but failed to check if limit equals function value. This forces Sloka to recognize that continuity requires ALL THREE conditions, not just one.

\vspace{0.5cm}


\subsection*{Problem 2: Day 3-4: The MVT Condition Bypass}

\begin{tcolorbox}[colback=blue!5!white, colframe=blue!50, arc=2mm]
  Let $f$ be a function with the following properties:
- $f(1) = 5$ and $f(4) = 14$
- $f$ is continuous on $[1, 4]$
- $f$ has a sharp corner at $x = 2.5$ (not differentiable there)

A student claims: "By the Mean Value Theorem, there exists $c$ in $(1, 4)$ such that $f'(c) = 3$ because $\frac{f(4) - f(1)}{4 - 1} = \frac{14 - 5}{3} = 3$."

**Your task:**
(a) Compute the average rate of change of $f$ on $[1, 4]$.
(b) Can the Mean Value Theorem be applied to $f$ on $[1, 4]$? Use the sentence frames to explain.
(c) Is the student's conclusion correct? Is the student's reasoning correct? Explain the difference.
\end{tcolorbox}

\textbf{The Trap:} TRAP: This is a CONDITION_BYPASS error. The student correctly computed the average rate but failed to verify differentiability on the open interval. The sharp corner at x = 2.5 means f is NOT differentiable on (1,4), so MVT doesn't apply. This problem forces Sloka to check BOTH conditions (continuity AND differentiability) before applying MVT. The twist: the conclusion might still be true (there might be a point where f'(c) = 3), but the reasoning is invalid because MVT can't be invoked.

\vspace{0.5cm}


\subsection*{Problem 3: Day 5-7: Integration Problem - Building to the FRQ}

\begin{tcolorbox}[colback=blue!5!white, colframe=blue!50, arc=2mm]
  Let $g$ be a function that is continuous on $[0, 8]$ and differentiable on $(0, 8)$. Selected values of $g$ are given in the table:

\begin{center}
\begin{tabular}{|c|c|c|c|c|}
\hline
$x$ & 0 & 3 & 5 & 8 \\
\hline
$g(x)$ & 2 & 11 & 8 & 14 \\
\hline
\end{tabular}
\end{center}

(a) Explain why there must be at least one value $c$ in $(0, 3)$ such that $g'(c) = 3$. Your explanation must include:
   - Verification of MVT conditions
   - Computation of average rate of change
   - Statement of conclusion

(b) A student claims that $g'(x) = 3$ for ALL $x$ in $(0, 3)$ because the average rate is constant. Is this reasoning valid? Explain.

(c) Is it possible that $g'(c) = 0$ for some $c$ in $(3, 5)$? Justify using appropriate theorems.
\end{tcolorbox}

\textbf{The Trap:} TRAP: This integrates multiple error types. Part (a) requires COMPLETE condition verification (targeting CONDITION_BYPASS). Part (b) tests whether student confuses average rate with instantaneous rate—a CER_BREAKDOWN error where evidence (average rate = 3) doesn't support the claim (derivative = 3 everywhere). Part (c) requires recognizing that g decreases from 11 to 8, so by MVT the average rate is negative, but Rolle's Theorem might also apply if we think about extrema. This builds complexity toward the main FRQ.

\vspace{0.5cm}


% ============================================================================
% FOOTER
% ============================================================================

\vfill

\begin{center}
  \small
  \textcolor{gray}{Generated by Alpha AP Justification Training System} \\
  \textcolor{gray}{11-03-2026} \\
  \textcolor{gray}{FRQ ID: undefined}
\end{center}

\end{document}
